\chapter{This kit's own deviations and ledger}
\label{app:ledger}

Chapter~\ref{ch:ledger} argues that an artefact should carry a dated list of
what it does not do. This appendix is that list, for the book you are holding.

\judgement{It is here for two reasons. The honest one: this book makes claims
and you should be able to see which of them it cannot support. The
demonstrative one: a chapter telling you to publish your gaps, in a repository
that publishes none, would be advice nobody should take.}

\section{The ledger}

Measured from the repository rather than typed. \code{tools/ledger.py} writes
these values and the build fails if they are stale, which is the same
discipline Chapter~\ref{ch:claims} asks of your application documents.

\begin{kittable}{@{}Xl@{}}
Parts & \val{parts} \\
Chapters & \val{chapters} \\
\quad of which bank chapters & \val{bankchapters} \\
Appendices & \val{appendices} \\
Bank questions & \val{questions} \\
Drills & \val{drills} \\
Words of chapter prose & \val{words} \\
\midrule
Reported claims & \val{reported} \\
\quad resting on one source, marked as such & \val{reportedsingle} \\
Statements marked as the author's judgement & \val{judgements} \\
\midrule
Sources & \val{sources} \\
\quad strong & \val{srcstrong} \\
\quad moderate & \val{srcmoderate} \\
\quad single-source & \val{srcsingle} \\
\quad recorded only to be refused & \val{srcrefused} \\
\midrule
Templates & \val{templates} \\
Role briefs & \val{rolebriefs} \\
Printable cards & \val{cards} \\
\end{kittable}

Two rows are worth reading against each other. There are \val{judgements}
marked judgements and \val{reported} sourced claims, which is roughly two
opinions for every reported fact.

\judgement{That ratio is not a defect and it is not hidden. This subject has
very little that can be measured, and a book on it that produced a citation for
every sentence would be citing content-marketing to do it \dash{}
Appendix~\ref{app:refused} is what that alternative looks like when examined.
What the ratio does mean is that a reader who disagrees with the author will
find a great deal to disagree with, and the marking is there so they can.}

\section{What this book has not done}

Each row is dated, carries a reason, and says what would close it \dash{}
because a list of gaps without closing conditions is a list of excuses.

\begin{kittable}{@{}lXX@{}}
\textbf{\#} & \textbf{What} & \textbf{What would close it} \\
\midrule
D-1 & No reader has used this to prepare for a process and reported back. The method's components are evidenced; the book as a whole is not. & Two or three people running a process with it and saying what was wrong \\
D-2 & Several primary sources were read through coverage rather than directly, because they were unreachable from the machine this was written on. \code{sources.json} marks each one. & Reading them, and re-checking what rests on them \\
D-3 & Appendix~\ref{app:reading} is seeded rather than complete & A pass whose whole purpose is the reading list \\
D-4 & Bank chapters \ref{ch:model} to \ref{ch:fundamentals} carry fewer questions and shorter contested sections than \ref{ch:retrieval} to \ref{ch:cost} & A second pass over Part~V, question by question \\
D-5 & The diagrams are few, and several chapters that would be clearer with one do not have one & A figure pass, with the widths measured rather than assumed \\
D-6 & No Polish edition, though the estate's other book has one. Bilingual is taken all the way \dash{} one body, a language file, a parity gate class \dash{} or not at all & Somebody wanting it enough to build the parity machinery \\
D-7 & Every page count and layout measurement was taken on a bare TeX installation with neither of the usual font packages. Line breaks, and therefore page counts, differ on a fuller one & Nothing. It is recorded so no page figure is quoted as a fact about both \\
\end{kittable}

\section{The one claim that stays open}

\textbf{This method is not validated end to end.}

Its components are evidenced to different degrees, and the book says which
throughout: the loop shape and the topic ordering rest on a corpus with a
stated method; the level definitions rest on published ladders; the evaluation
practice rests on practitioner writing. The assembly of those into a method,
and the claim that following it improves anything, rests on nothing measured.

\judgement{That is the same standing as the 80/80 claim in the estate's other
book, and it is recorded the same way: printed, dated, and not quietly dropped.
It will stop being true when several people have used this and said what went
wrong, and until then it is the first thing an honest reader should hold
against it.}

\section{Where the register lives}

\code{DEVIATIONS.md}, in the repository, with dates. This appendix is printed
from the same reasoning but it is not generated from that file \dash{} which is
itself a deviation, of exactly the kind Chapter~\ref{ch:ledger} warns about,
and it is D-8 in the file rather than being left for a reader to notice.
